\section{Rust Runtime}

Starting from the demo application provided by the \textit{Ada} guide
\cite{burnsGuideUseAda2004}, we then went on implementing an equivalent system
fully written in \textit{Rust}. 

Compared to \textit{Ada}, \textit{Rust} does not provide a reference runtime
implementation, leaving to third-party developers this incumbency. This is
advantageous from a flexibility and ecosystem development standpoint, since it
encourages the growth of a competitive runtime ecosystem, and fosters the
development of multiple different implementation, tailored to more specific
use-cases. On the other hand, this clearly implies more ecosystem fragmentation,
and could also result in lower "quality" of the runtimes, since a single
reference implementation will naturally tend to concentrate resources for its
development from all the actors interested in its usage or analysis, instead of
dispersing them across multiple projects. 

Thus, when designing and implementing a real-time system in \textit{Rust}, some
effort must be spent in order to identify and integrate the most appropriate
runtime components for the system, having care to select compatible components
that can work in harmony together. In our case this means choosing a
\textit{concurrency runtime} that provides the task abstraction and scheduling
primitives, a \textit{logging framework} and a \textit{Hardware Abstraction
Layer} that provides higher-level abstractions over the system hardware. 

\subsection{RTIC}

As for the \textit{concurrency runtime}, we have decided to use the
\textit{RTIC} framework \cite{RticrsRtic2025}. It is generally considered to be
a "concurrency framework", since it does not contain a full-fledged software
kernel and relies on external \textit{Hardware Abstraction Layer} to access
peripherals. 

\textit{RTIC} implements the \textit{Stack Resource Policy} for shared resource
access management \cite{bakerStackbasedSchedulingRealtime1991}, with
\textit{Fixed-Priority Scheduling} \cite[Chapter~6]{liuRealTimeSystems2000}. 

\textit{SRP} has some really useful runtime properties for real-time systems,
both for performance and especially for amenability to response-time analysis.
Some of those properties are: 
\begin{itemize}
  \item Preemptive deadlock and race-free scheduling;
  \item Tasks execute on a single shared stack;
  \item Tasks run-to-completion with wait free access to shared resources;
  \item Bounded priority inversion by a single critical section.
\end{itemize}
Those come with the following set of requirements:
\begin{itemize}
  \item Tasks are associated with static \textit{pre-emption levels};
  \item Tasks execute on a single core;
  \item Tasks do not voluntarily suspend (run-to-completion);
  \item Resource must be claimed and released in LIFO order.
\end{itemize}
The motives behind these requirements and the proof of the properties are
outside the scope of this work; they can be found in the original paper
\cite{bakerStackbasedSchedulingRealtime1991}. 

\textit{SRP} based scheduling requires the set of static priority tasks and
their access to shared resources to be known in order to compute a static
ceiling $\pi$ for each resource. The static resource ceiling $\pi(r)$ reflects
the maximum static priority of any task that accesses the resource $r$. This is
the reason why, despite its excellent properties, is not used in general-purpose
OS: it requires the resource dependency graph to be known statically, which is
generally not possible. 

\textit{FPS} is a very simple scheduling policy: every task is statically given
a fixed priority value (values must have an order), and all jobs released will
inherit the priority of their task. At every point in time the running job is
the one with the highest priority among all the \textit{ready} ones. 

The key idea of \textit{RTIC} is to use the native features provided by the
interrupt engine of ARM microcontrollers (\textit{NVIC}
\cite{armlimitedARMCortexM42009}) to implement these two protocols, obtaining a
\textit{hardware-accelerated runtime} suitable for real-time systems, with
minimal overhead.  

In the case of the ARM Cortex-M architecture, each interrupt vector entry $v_i$
is associated an \textit{Interrupt Service Routine} $isr_i$, a static priority
($p_i$), an enabled- ($en_i$) and a pending-bit ($pend_i$). 

An interrupt $i$ is scheduled (run) by the hardware only if both of the
following conditions are met:
\begin{enumerate}
  \item $i$ is pending and enabled and has a priority higher than the
  \texttt{BASEPRI} register;
  \item $i$ has the highest priority among interrupts meeting 1.
\end{enumerate}

The \textit{SRP} model for single-core static scheduling, on the other hand,
states that a task should be scheduled (run) under the conditions: 
\begin{enumerate}
  \item It is requested to run and has a static priority higher than the current
  system ceiling ($\Pi$); 
  \item It has the highest static priority among tasks meeting 1.
\end{enumerate}
The system ceiling $\Pi$ is computed as the maximum priority ceiling of the
currently held resources. 

This shows a very close relationship between the scheduling logic implemented in
hardware by the NVIC, and the model required by \textit{SRP}. The mapping
between logical tasks and actual hardware is quite straightforward: 
\begin{itemize}
  \item Each task $t$ is mapped to an interrupt vector index $i$ with the
  corresponding ISR $isr_i = t$ and given the correct static priority $p_i =
  p(t)$; 
  \item The system priority ceiling $\Pi$ is mapped to the \texttt{BASEPRI}
  register, which is updated at the entry and exit of every \textit{critical
  section} (i.e. lock). 
\end{itemize}
The mapping between tasks and ISR, and resources to their static ceilings is
computed at compile-time by \textit{RTIC} using \textit{procedural macros}, thus
is has no impact on the performance at runtime. 

Actually, \textit{RTIC} maps each task priority level to an ISR at the
corresponding priority. Each ISR runs an \textit{asynchronous executor} that
handles tasks wrapped as \textit{async} functions assigned to that priority
level. This allows for \textit{async/await} style programming in tasks, lifting
the run-to-completion requirement for them, since each section of code between
two yield points can be seen as an individual task. The compiler will reject any
attempt to await while holding a resource (not doing so would break the strict
LIFO requirement on resource usage under SRP) \cite{RticrsRtic2025}. This in
turn gives much better developer ergonomics, especially for managing inter-task
communication and for handling interaction with the hardware. 

\subsection{Defmt}

The second piece of the system is \textit{defmt} \cite{IntroductionDefmtBook}:
it is a logging framework for embedded applications, written in \textit{Rust}. 

It uses a novel approach to improve space and computational efficiency of the
logging on the embedded device side. Instead of storing the format strings and
perform parameter formatting on the device, it \textit{defers} those operations,
moving them on the connected host. 

Technically, it embeds metadata about format string and the parameters types on
a linker section that gets stripped before placing the firmware into the device
flash. Thus, the device only needs to send the format string index and raw
parameters; then the connected host, provided with the non-stripped binary file,
will perform formatting and printing. 

The disadvantage of this approach is that without the metadata embedded in the
binary file, the messages from the device are not interpretable. Nonetheless,
the advantages in terms of size and computational requirements outweighs the
downsides in most cases. 

Defmt supports different \textit{backends} for communication between the host
and the device, such as \textit{RTT}, \textit{semihosting} and \textit{ITM}; in
our implementation we opted to use \textit{RTT}, as we had problems with the
semihosting backend. 

\subsection{HAL}

The final key component of the system is the \textit{Hardware Abstraction Layer
(HAL)}. It provides a high-level, portable interface to interact with the
underlying hardware peripherals, abstracting away the low-level register
manipulation and platform-specific details. 

In the \textit{Rust} embedded ecosystem, \textit{HAL}s are typically implemented
as crates that follow the \textit{embedded-hal} trait specifications
\cite{RustembeddedEmbeddedhal2025}. This standardization allows for
interoperability between different hardware platforms and peripheral drivers,
enabling code reuse and ecosystem compatibility. 

For our implementation, we used the \textit{STM32F4xx HAL} crate
\cite{Stm32rsStm32f4xxhal2025}, which provides abstractions for \textit{GPIO}
pins, timers, \textit{UART} communication, \textit{SPI}, \textit{I2C}, and other
peripherals specific to the \textit{STM32F4} microcontroller family. The
\textit{HAL} ensures type-safe access to hardware resources at compile time,
preventing common errors such as accessing uninitialized peripherals or
conflicting resource usage, enforced by \textit{Rust} strong type system. 
In general \textit{RTIC} is \textit{HAL} agnostic: it is possible to use other
implementations from different vendors (e.g. \textit{Embassy}). 

The most important peripheral needed for the experiment is the timer, which is
used to provide the notion of time to the system, via a global
\textit{rtic-monotonics} object, thus enabling time-based operations (e.g.
delays, periodic tasks, deadline monitoring). \textit{RTIC} has built-in support
for our chosen \textit{HAL}, hence implementing the required timer object only
required calling a simple macro \cite{Rticmonotonics}.  